% Short running form; the full title is too wide for the head's right box.
\runninghead{The Song of the Exiles}
\properlane{songs-of-sion-in-babylon}{The Song of the Exiles: Confession, Humility and Hope by the Rivers of Babylon}
\label{sec:songs-of-sion}

Two of the formulary's chants are songs of exile. The Introit is cut from the
prayer Azarias prayed in the Babylonian furnace; the Offertory is the first
verse of the captives' lament by Babylon's rivers. The Fathers read both of the
present condition of God's people. St Augustine teaches that the citizens of
Jerusalem are captives in a world that flows past like Babylon's rivers, and
that they are to sit by those rivers in humility, not plunge into them, and to
weep in the memory of Sion; St Hilary of Poitiers, that the psalm speaks of a
spiritual captivity, the soul's exile since Adam from the heavenly Jerusalem,
its mother. St Jerome teaches the exiles' first word, to confess in any distress
that whatever befalls is borne justly, and St Augustine places that confession
first in the order of prayer, before any petition. The Introit's verse names
the road home, the law of the Lord, and the Communion names the exiles' hope,
God's word, which comforts in humiliation. Read with these, the formulary sets
an exiled people's confession and lament beside an Epistle that tells the same
people, in evil days, to redeem the time and to sing in their hearts,
\latin{cantantes, et psallentes}, and an Alleluia whose ready heart answers in
the same two verbs, \latin{cantabo, et psallam}. The captives of the psalm ask,
\latin{Quomodo cantabimus canticum Domini in terra aliena?} The Mass answers
with a song of the heart that confesses God's justice, gives thanks in
adversity, and waits for the city it remembers.

\subsection{The confession in the furnace: St Jerome and St Augustine}

Azarias's prayer begins with praise of God's justice and goes on to confession:
``For thou art just in all that thou hast done to us \dots\ For we have sinned,
and committed iniquity'' (Dan 3:27, 29). St Jerome reads the opening verse as a
rule for every Christian in trouble: \latin{Quando diversis premimur angustiis,
ex toto cordis hoc loquamur affectu, et quidquid nobis acciderit, juste nos
sustinere fateamur}, when we are pressed by distresses of every kind, let us say
this with the whole affection of the heart and confess that whatever has
befallen us we bear justly. On the confession of sin he makes a harder point.
\latin{Et certe tres pueri non peccaverant}: the three youths had not sinned,
and were not of an age to be punished for their own faults when they were
carried to Babylon; \latin{ergo \dots\ ex persona populi loquuntur}, they speak
in the person of the people. And on the verse that mourns the loss of prince,
prophet and sacrifice (3:37--38), Jerome gives the Church leave to pray it as
her own: these verses are to be used \latin{si quando Ecclesiae propter peccata
populi \dots\ sustinent penuriam, et quando in persecutionibus non offertur
sacrificium et oblatio}, whenever the Churches, for the sins of the people,
suffer want of holy teachers, and when in persecution the sacrifice is not
offered. Of the angel who goes down into the furnace (3:49) he says that it is
the divine word, \latin{sermo divinus}, which comes to the soul that has
despaired of human help and turned wholly to the Lord (\work{Commentarii in
Danielem}, on 3:26, 29, 37, 49).

St Augustine finds the same confession at the head of all the saints' prayer.
Expounding the Gradual's own psalm, he comes to ``The Lord is just in all his
ways'' (Ps 144:17), and says: \latin{omnes sancti in tribulationibus
constituti, prius justitiam ejus laudaverunt, et sic beneficia petiverunt. Prius
dixerunt: Justum est quod facis. Sic rogavit Daniel}, all the saints, set in
tribulation, first praised his justice and so asked his benefits; first they
said, What you do is just; so Daniel asked (\work{Enarrationes in Psalmos} 144,
on v.~17). ``First they praised Him scourging,'' the English renders his close,
``and so they felt Him feeding.'' The Introit's order is that order: the
confession that God's judgments are true, then the plea, \latin{sed da gloriam
nomini tuo}.

\subsection{The road home: the Introit's verse}

The verse after the antiphon, ``Blessed are the undefiled in the way, who walk
in the law of the Lord'', names the way back. St Augustine opens his long
exposition of Ps~118 by observing that the psalm begins with blessedness, which
no one does not want: \latin{Scio quid velis, beatitudinem quaeris; si ergo vis
esse beatus, esto immaculatus. Illud enim omnes, hoc autem pauci volunt, sine
quo non pervenitur ad illud quod omnes volunt. Sed ubi erit quisque
immaculatus, nisi in via? In qua via, nisi in lege Domini?} I know what you
want, he makes the psalm say, you seek blessedness; if then you would be
blessed, be undefiled; all want the one, few the other, without which the one
all want is not reached; and where will anyone be undefiled but in the way, and
in what way but in the law of the Lord? The lustful, the avaricious, the cruel
and the ambitious all seek blessedness, by evil roads, and the verse is the
divine voice calling them back: \latin{Quo itis? Peritis, et nescitis}, where
are you going? You are perishing and do not know it (\work{Enarrationes in
Psalmos} 118, sermo 1, 1).

St Hilary reads the same verse as the first stage of an ascent. The order of the
sayings matters: first ``Blessed are the undefiled in the way'', then ``Blessed
are they that search his testimonies'', because a man first enters the way of
truth with his morals made firm, and only then searches God's testimonies with a
cleansed mind. The undefiled are those who curb the flesh, tame the mind's
wantonness, conquer the hunger of avarice and shun the glory of earthly honours;
and the words \latin{in lege Domini} are not idle, for there is another law, the
law of sin in the members (Rom 7:23; \work{Tractatus super Psalmos}, in
Ps.~118, \latin{Aleph} 1--3). St Robert Bellarmine joins them: the law of the
Lord is a straight and clear path, ``opposed to the law of the flesh, which the
Apostle designates as `the way of concupiscence,' full of the dust of pride, the
mud of luxury, and the dirty water of avarice''; and ``Christ himself informs us
that the straight road to eternal happiness is the observance of his law''.
For the exiles of the Introit, then, the verse names the road out of Babylon:
not a journey on foot but the keeping of the commandments they confess they
broke. The Postcommunion asks for that keeping, \latin{tuis semper oboedire
mandatis}, and the Epistle begins with the same verb as the verse, \latin{Videte
quomodo caute ambuletis}, walk carefully, \latin{non quasi insipientes, sed ut
sapientes}, as Augustine's undefiled walk the way and not the evil roads.

\subsection{By the waters, not in them: St Augustine and St Hilary on the Offertory}

St Augustine preached on the Offertory's psalm on a day when it had been sung,
and his sermon is a programme for the exiled Church. ``The waters of Babylon are all
things which here are loved, and pass away.'' One man loves farming, another
soldiering, another office; let each see ``that what he hath loved is not a
foundation of Jerusalem, but a stream of Babylon''. The citizens of the holy
Jerusalem, ``understanding their captivity'', see these waters dragging men
down to the sea, and ``they throw not themselves into the waters of Babylon, but
`sit down and weep,' \dots\ sitting, that is, humbling themselves''. ``Sit `by'
the waters, not `in' the waters, not `under' the waters; but yet sit, in humble
fashion.'' Their weeping is not Babylon's: ``Thou oughtest to weep, but in the
remembrance of Sion. If thou weepest in the remembrance of Sion, thou oughtest
to weep even when it is well with thee in Babylon.'' Their instruments are
``God's Scriptures, God's commands, God's promises, meditation on the life to
come'', and they hang them on the willows, barren trees watered by Babylon,
because such men are unfit to carry them (\work{Enarrationes in Psalmos} 136, on
vv.~1--2). He begins the sermon by saying that the song must be done as well as
sung, \latin{oportet ut non tantum ista decantemus, sed et faciamus}, and he
ends it, \latin{Fratres, organa in operando non cessent; vobis invicem cantate
cantica Sion}: brothers, let not your instruments rest in your work; sing to
one another the songs of Sion.

St Hilary reads the same psalm more inwardly still. Its captivity is spiritual,
not bodily; Sion is the seat of eternal blessedness and \latin{mater coelestis
nobis Jerusalem}, our heavenly mother Jerusalem; and man has been an exile from
her since Adam's sin. He asks those who read every psalm literally which rivers
and which willows could be meant, and he answers that the rivers are the works
of the age and of the body, which flow past. The willows, which revive in water,
are the saints and faithful, revived by the word of God and the water of baptism
(Isa 44:4); the instruments are our bodies, by whose honest movements we sing
what pleases God, hung up in continence; and the saint will not share the songs
of Sion with the powers that hold him captive (\work{Tractatus super Psalmos},
in Ps.~136). Here Augustine and Hilary disagree in a shared reading. For
Augustine and Bellarmine the willows are barren men; for Hilary, and Cassiodorus
after him, they are the saints. For Augustine the instruments are the
Scriptures, for Hilary our bodies; Cassiodorus takes Augustine's instruments and
Hilary's willows, and has the Scriptures hung on holy men when grace is shared.
Cassiodorus draws a posture from the psalm's first words, \latin{super
flumina}: those already freed sit on the banks, not in the rivers, and sit in
humility; and he calls the singing of psalms the consolation of those who weep,
the care of those who grieve and the health of the sick, a remedy of the soul (\work{Expositio Psalmorum} 136.1).

St John Chrysostom reads the psalm of the historical captives, and draws from
them a warning for the Church (\work{Expositio in Psalmos}, on Ps~136). While
they held the city they scorned it, and God cast them out to teach them by loss
to desire it; those who had mocked the weeping prophets now wept unbidden and
kept the law among their captors. So, he says, many will seek the Jerusalem
above on the last day and find no return; the present life must be ordered now,
so as not to become strangers to that city.

\subsection{Singing in evil days: the Epistle and the Alleluia}

The Epistle speaks to the same people in the same condition. ``Redeeming the
time, because the days are evil'': St Jerome explains that the wise man buys
back the time that human malice had sold, \latin{dies malos in bonos vertimus},
we turn evil days into good, as Joseph and Job were changed by no change of
times. ``Singing and making melody in your hearts'': \latin{Audiant haec
adolescentuli: audiant hi quibus psallendi in ecclesia officium est, Deo non
voce, sed corde cantandum}, let the young men hear it, let those whose office it
is to sing in church hear it, that God is to be sung to not with the voice but
with the heart; a singer of poor voice whose works are good \latin{dulcis apud
Deum cantor est}, is a sweet singer before God. ``Giving thanks always for all
things'': if a house falls, if a beloved wife and children are taken into
captivity, the Christian's word is \latin{Benedictus Deus, minora me scio
sustinere quam mereor}, blessed be God, I know that I suffer less than I deserve
(\work{Commentarii in Epistolam ad Ephesios} III, on 5:16, 19, 20). St John
Chrysostom gives the same verses the exile's turn: ``The time is not yours. At
present ye are strangers, and sojourners, and foreigners, and aliens''; redeem
it by giving up whatever is demanded, ``only preserve the principal thing, I
mean the faith''. Against drunkenness he sets the psalms: ``Learn to sing
psalms, and thou shalt see the delightfulness of the employment. For they who
sing psalms are filled with the Holy Spirit''; and ``with your hearts'' means
``with close attention and understanding''. Thanks are owed in everything, ``be
it even disease, be it even penury'' (\work{Homilies on Ephesians} 19). Jerome
subordinates the voice to the heart and Chrysostom asks for the voice with
attention; both oppose singing that seeks only to please, and Bellarmine,
expounding the Offertory's psalm, makes the same complaint of those who ``sing
the song of the Lord in a strange land'' when they sing sacred hymns only to
please the ear.

The Alleluia answers in the Epistle's own two verbs: \latin{cantantes, et
psallentes} there, \latin{cantabo, et psallam} here. Its verse is also
Ps~56:8, and St Augustine reads it of a heart made ready by tribulation: ``the
patience of good men with preparation of heart accepteth the will of God: and
glorieth in tribulations''. His example is St Paul, in chains, in prisons and in
hunger, ``saying, `We glory in tribulations.' Whence, but that prepared was his
heart? Therefore he was singing and playing'' (\work{Enarrationes in Psalmos}
56, on v.~8). The captives of the Offertory could not sing the Lord's song in a
strange land; the prepared heart of the Alleluia sings it, not to the captors
but to God, \latin{psallam tibi}.

\subsection{Fed in due season, comforted in humiliation}

The other elements belong to the same exile. The Gradual lifts the eyes of all
to the God who gives food in due season, and Augustine
reads the delay as a physician's care; the exiles are fed, and they wait. The Communion asks
God to remember the word in which he gave hope, \latin{haec me consolata est in
humilitate mea}. Augustine reads the humiliation as tribulation, whether pride
deserved it or patience proves it, and says: ``This hope was given to Christ's
Body, that is, to the Church, that it might be a comfort to Her in her
humiliation'' (\work{Enarrationes in Psalmos} 118, sermo 15, on v.~50). St
Hilary hears in the verse the consolation that God's every word gives in
sickness, persecution, loss and bereavement, since every word of God calls us
to the hope of heavenly goods (\work{Tractatus super Psalmos}, in Ps.~118,
\latin{Zain}). Azarias's prayer has the same word: \latin{sumusque humiles in
universa terra hodie propter peccata nostra}, we are brought low in all the
earth this day for our sins (Dan 3:37).

The orations pray as exiles who have confessed. The Collect asks pardon and
peace. Sicard of Cremona and William Durandus, expounding these chants of captivity,
say that the bodily captivity signifies our spiritual one, and that the
return from captivity is the remission of sins (\work{Mitrale} VIII.20;
\work{Rationale} VI.137). The Secret asks heavenly medicine for the heart's
vices, as Cassiodorus called the psalmody of the exiles the soul's remedy. And
the Postcommunion asks the obedience by which, as the Introit's verse taught,
the exiles walk home. At the offertory of every Mass, too, the priest prays in
Azarias's words, \latin{In spiritu humilitatis et in animo contrito suscipiamur
a te, Domine}: the furnace prayer offers the sacrifice of a people that has no
other.

The Gospel enters the reading as one exile's crisis. A son lies dying in a world
that passes. St Gregory preached it at a martyrs' tomb in a time when, he says,
\latin{Ubique mors, ubique luctus, ubique desolatio}, death, grief and
desolation were everywhere; the world that once held men by delight now
\latin{tantis plagis plenus est, ut ipse nos jam mundus mittat ad Deum}, is so
full of blows that it sends us to God itself (\work{Homiliae in Evangelia}
28, 3). The father's journey to Christ is that sending, and the word that reaches
his son is the word Jerome finds in the furnace, the angel who is \latin{sermo
divinus}.

\subsection{Two hard sayings}

The first is in the Introit. Read of every suffering, ``all that thou hast done
to us \dots\ thou hast done in true judgment \dots\ for we have sinned'' could be
taken to say that all suffering is punishment for the sufferer's own sin. St
Jerome closes that reading: the youths who prayed it had not sinned, and spoke
\latin{ex persona populi}. Augustine places the confession in the order of
prayer, praise of God's justice before petition, and not in an account of the
causes of every pain. The confession belongs to the people praying, and it may
be prayed by the innocent for them.

The second is the close of the Offertory's psalm, which the Mass does not sing:
``Blessed be he that shall take and dash thy little ones against the rock''
(Ps 136:9). Five witnesses read it. St Augustine asks, ``What are the little
ones of Babylon? Evil desires at their birth'': ``when lust is little, \dots\
dash it. But thou fearest, lest though dashed it die not; `Dash it against the
Rock; and that Rock is Christ'.'' St Hilary reads the daughter of Babylon as the
flesh and her little ones as \latin{tenera adhuc corporis vitia}, the body's
vices while still tender, to be crushed against the rock, since \latin{Petra
autem, secundum Apostolum, Christus est} (\work{Tractatus super Psalmos}, in
Ps.~136, 14). Cassiodorus takes the little ones as the flesh's errors while
they are small, to be dashed at once on the rock, the Lord Saviour. St Robert
Bellarmine gives first the literal prophecy, that the Medes and Persians would
dash Babylon's children, and then a spiritual sense in which the little ones
are persons: Babylon's young, converted and brought ``to the rock, Christ, to
get a fortunate dash against it, and die the death of the old man, to rise a
new man'', and, morally, those not yet deeply stained by sin, for ``blessed is
he, who, by a happy dash on the rock, kills sin''. St John Chrysostom reads no
figure in the verse at all: the words are the captives' own anger, which the
prophet reports and does not make his own, and the New Covenant commands
otherwise, that we feed our enemies and pray for those who abuse us
(\work{Expositio in Psalmos}, on Ps~136:8--9). Augustine, Hilary and Cassiodorus
make the little ones the beginnings of sin; Bellarmine makes them persons
brought to Christ to die to sin; all four make the rock Christ. Chrysostom
alone leaves the verse as the captives' passion, and answers it with the
Gospel.

\subsection{The four senses of this reading}

\begin{fourSenses}
\item[Literal.] In the Babylonian furnace Azarias confesses that God's
judgments on his people were just and asks mercy for his name's sake (Dan
3:26--45); captives by Babylon's rivers weep in the memory of Sion and refuse
to sing its songs for their captors (Ps 136); St Paul tells the Ephesians to
walk wisely, to redeem the time in evil days, to be filled with the Spirit and
to sing in their hearts, giving thanks always (Eph 5:15--21).

\item[Allegorical.] Babylon is the world whose loves flow past, Sion--Jerusalem
the heavenly city and mother of the faithful, and the captives are the Church's
citizens on pilgrimage (Augustine; Hilary); the instruments hung up are the
Scriptures (Augustine) or the body kept in continence (Hilary); the rock
against which Babylon's little ones are dashed is Christ (Augustine; Hilary);
the angel in the furnace is the divine word (Jerome).

\item[Moral.] Confess God's justice before asking his gifts (Jerome;
Augustine); walk the way of the Lord's law and not the roads to a false
blessedness (Augustine; Hilary); sit by the rivers and not in them, and weep in
the memory of Sion even when it is well with you here (Augustine); kill sin
while it is small, against Christ (Augustine; Hilary); redeem the time by good
works, sing with the heart and give thanks in loss (Jerome; Chrysostom).

\item[Anagogical.] The memory of Sion is longing for the heavenly Jerusalem:
``whither your hope goeth before, let your life follow \dots\ Your captivity
will pass away, your happiness will come; the last enemy shall be destroyed,
and we shall triumph with our King, without death'' (Augustine, on Ps~136); the
hope of the Communion is for the promises kept in glory, to which God's every
word calls (Hilary).
\end{fourSenses}
